% Frames cut from session2_slides.tex on 2026-10-26. They belong to Session 3 of
% admin/course_plan_4_sessions.md (blocks 23-25). Not compiled on its own: move them
% into the Session 3 deck, then delete this file. Frames are verbatim, grouped by
% destination block.
%
% Edits still needed when they land in Session 3:
%   - Titles: Part 2 says "(30 min)"; block 24 is 25 min. Part 3 says "(30 min)"; block 25
%     (robustness workflow plus Part 3) is 25 min.
%   - "Today's data" / "today" wording assumes the data were introduced in the same
%     session; in Session 3 participants reopen their Session 2 folder.
%   - "What to Watch For" and "Debrief" still ask about the merge (Part 1, now Session 2).
%   - "Two More Safeguards" (second agent as critic) overlaps the Session 2 checker tab.

% ══════════════════════════════════════════════════════════════════
%  Block 23: verification extras
% ══════════════════════════════════════════════════════════════════

\begin{frame}{The Requirements-to-Tests Matrix}
  \footnotesize
  \renewcommand{\arraystretch}{1.2}
  \begin{tabular}{@{}L{0.27\textwidth}L{0.21\textwidth}L{0.42\textwidth}@{}}
    \toprule
    \textbf{Requirement} & \textbf{Test} & \textbf{Passing condition} \\
    \midrule
    Merge keeps every respondent & Row-count and key test & 2,044 rows; \texttt{respondent\_id} unique \\
    Outcome built as defined & Variable-construction test & \texttt{supports\_intervention} is 1 exactly for ``some'' or ``big'' \\
    Control is the reference arm & Specification test & The omitted category is \texttt{cntrl} \\
    Survey weights are used & Independent recomputation & Logistic of the base-model intercept equals the weighted control share from the raw CSV \\
    N is explained for every model & Sample accounting & Rows in $-$ rows dropped (by reason) = N used \\
    Every specification is reported & Output manifest & All specifications are in the table, each with its N \\
    \bottomrule
  \end{tabular}

  \vspace{0.5em}
  \small
  \alert{Write the matrix before the code.} A requirement with no test is unverified.
\end{frame}

\begin{frame}[fragile]{The Audit Report: Evidence, Not ``All Tests Passed''}
  \begin{lstlisting}[basicstyle=\ttfamily\footnotesize]
## Audit report
Tests run:         commands; passed / failed / skipped
Requirements:      each requirement -> its test
Not tested:        anything without a test, and why
Numerical checks:  our value vs. independent value
Needs your call:   choices Claude proposed
Not covered:       identification, construct validity
  \end{lstlisting}

  \vspace{0.3em}
  \begin{block}{What the report is for}
    Not for Claude Code to declare the research correct. It documents what was exercised, against which expectation, and what remains uncertain.
  \end{block}
\end{frame}

\begin{frame}{Two More Safeguards}
  \begin{columns}[T]
    \begin{column}{0.48\textwidth}
      \textbf{Replicate before you extend}
      \begin{enumerate}
        \small
        \item Give Claude Code the original code or description
        \item Ask it to replicate a key table
        \item Compare with the published result
        \item Only then begin your extension
      \end{enumerate}
      {\small If you cannot reproduce the baseline, nothing built on it is interpretable.}
    \end{column}
    \begin{column}{0.48\textwidth}
      \textbf{Use a second agent as a critic}
      \begin{enumerate}
        \small
        \item Agent 1 writes the code
        \item Open a new conversation (Agent 2) and ask: ``Read the analysis script and identify any bugs, silent data loss, or decisions that could change the results.''
      \end{enumerate}
      {\small Agent 2 has fresh context and no attachment to the code.}
    \end{column}
  \end{columns}

  \vspace{0.8em}
  \alert{Treat your agents like co-authors who check each other's work, not a single assistant you trust unconditionally.}
\end{frame}

\begin{frame}{What to Delegate, What to Verify, What to Keep}
  \begin{columns}[T]
    \begin{column}{0.32\textwidth}
      {\color{green!50!black}\textbf{Delegate freely}}\\[0.3em]
      {\small
      Boilerplate \& docs\\
      File management\textsuperscript{\dag}\\
      Reformatting code}
    \end{column}
    \begin{column}{0.32\textwidth}
      {\color{accent}\textbf{Delegate + verify}}\\[0.3em]
      {\small
      Data loading \& cleaning\\
      Data merges\\
      Variable construction\\
      Summary statistics\\
      Robustness variants\\
      Visualization\\
      \LaTeX{} tables}
    \end{column}
    \begin{column}{0.32\textwidth}
      {\color{red!70!black}\textbf{Do not delegate}}\\[0.3em]
      {\small
      Choosing specifications\\
      Interpreting results\\
      Identification strategy\\
      Sample selection\\
      Causal claims}
    \end{column}
  \end{columns}

  \vspace{0.5em}
  {\scriptsize \textsuperscript{\dag} Be careful with file management --- Claude Code can delete files. Always review before approving.}

  \vspace{0.5em}
  The test: \textit{If a referee questioned this choice, who would need to answer --- you or the RA?}

  \vspace{0.5em}
  Today's exercise focuses on the middle column: tasks where Claude Code is productive but where you must verify every output.
\end{frame}

% ══════════════════════════════════════════════════════════════════
%  Block 24: Exercise Part 2
% ══════════════════════════════════════════════════════════════════

\begin{frame}{Exercise Part 2: Specify and Estimate (30 min)}
  \begin{enumerate}
    \setcounter{enumi}{3}
    \item \textbf{Base model:} Logistic regression of \texttt{supports\_intervention} on \texttt{treatment}.
    \begin{itemize}
      \small
      \item Use survey weights. What should the reference category be? Tell Claude Code explicitly.
      \item \alert{Independent check:} the logistic of the intercept should equal the weighted control-group share, computed straight from the raw CSV
    \end{itemize}

    \vspace{0.2em}
    \item \textbf{Add FX exposure.} Did the sample size change? Why or why not?

    \vspace{0.2em}
    \item \textbf{Add demographic controls:} age, female, education, urban/rural, income, left-right.
    \begin{itemize}
      \small
      \item \alert{Before running: predict how many observations you will have.} Then check.
    \end{itemize}

    \vspace{0.2em}
    \item \textbf{Add an interaction} between treatment and FX exposure.
  \end{enumerate}

  \vspace{0.4em}
  {\small \alert{For every model, reconcile N:} rows in $-$ rows dropped (by reason) = N used.}
\end{frame}

\begin{frame}{Part 2 Debrief: What Happened to Your Sample?}
  \textbf{How many observations did you have in each model?}

  \vspace{1em}
  \pause
  \begin{alertblock}{The Missing Data Trap}
    Adding control variables can silently change your sample.
  \end{alertblock}

  \vspace{0.3em}
  \begin{tabular}{lrl}
    \toprule
    \textbf{Specification} & \textbf{N} & \textbf{What happened} \\
    \midrule
    Treatment only & 2,036 & Full sample (8 missing outcome) \\
    + FX exposure & 2,035 & 2 missing FX status \\
    + age, female, education, urban/rural & 2,035 & No additional missingness \\
    + income quintile & $\sim$1,473 & \color{red!70!black}{565 missing income} \\
    + left-right scale & $\sim$1,499 & \color{red!70!black}{542 missing left-right} \\
    + both & $\sim$1,151 & \color{red!70!black}{44\% of sample dropped} \\
    \bottomrule
  \end{tabular}

  \vspace{0.5em}
  \textbf{Did Claude Code mention this when it added the controls?}
\end{frame}

% ══════════════════════════════════════════════════════════════════
%  Block 25: robustness workflow, then Exercise Part 3 (plan order: workflow first)
% ══════════════════════════════════════════════════════════════════

\begin{frame}{The Robustness Workflow}
  \centering
  {\large You specify the core analysis.\\
  Claude Code generates the variations.\\
  You verify each one.}

  \vspace{1.5em}
  \textbf{Claude Code never chooses what robustness checks to run.}

  \vspace{1em}
  \textbf{For each specification, check:}
  \begin{itemize}
    \item Does it use the right sample? (What is N?)
    \item Does the code match what you asked for?
    \item Are the coefficients in a plausible range?
    \item Do the signs on control variables make sense?
  \end{itemize}

\end{frame}

\begin{frame}{Exercise Part 3: Robustness and Output (30 min)}
  \begin{enumerate}
    \setcounter{enumi}{7}
    \item \textbf{Generate robustness checks} from your base model:
    \begin{itemize}
      \item With and without survey weights
      \item Different control sets (track how N changes each time!)
      \item Collapse three treatment arms into ``any information'' vs.\ control
      \item Restrict to complete cases so N is constant across specs
    \end{itemize}

    \vspace{0.3em}
    \item \textbf{Format as \LaTeX.} Produce a publication-style regression table with all specifications side by side, and compile to PDF.
    \begin{itemize}
      \small
      \item Output manifest test: every specification appears in the table, each with its N
    \end{itemize}

    \vspace{0.3em}
    \item \textbf{Audit report.} Ask Claude Code for \texttt{output/audit\_report.md}, in the format from the audit-report slide.

    \vspace{0.3em}
    \item \textbf{Commit.} Save all scripts, tests, output, and the \LaTeX{} table.
  \end{enumerate}
\end{frame}

\begin{frame}{\LaTeX{} Tables: Bringing Results Into Your Paper}
  Until now, your results are printed to the console or saved as CSVs.\\
  In your actual workflow, they need to end up in a paper.

  \vspace{0.8em}
  \textbf{Why this matters:}
  \begin{itemize}
    \item Claude Code can write \LaTeX{} tables directly from regression output
    \item The table lives in your project directory, under version control
    \item When you update the analysis, Claude Code updates the table
    \item No manual reformatting, no copy-paste errors
  \end{itemize}

  \vspace{0.8em}
  \textbf{Today:} generate 5 tables instead of one table. In Session 3, we bring the full paper into VS Code.

  \vspace{0.5em}
  \begin{block}{The bigger idea}
    The more of your research that lives inside the project directory, the more Claude Code
    can operate on. Code, data, output, tables --- and eventually the paper itself.
  \end{block}
\end{frame}

% ══════════════════════════════════════════════════════════════════
%  End of the Swiss franc exercise (old Session 2 closing frames)
% ══════════════════════════════════════════════════════════════════

\begin{frame}{What to Watch For During the Exercise}
  This exercise is designed to surface moments where Claude Code makes choices for you.

  \vspace{0.5em}
  \begin{itemize}
    \small
    \item \textbf{The merge:} Did Claude Code pick the right join type? Did it mention the extra rows?
    \item \textbf{Reference categories:} Did it set treatment to ``control'' or pick alphabetically?
    \item \textbf{Missing data:} When you added controls, how many observations did you lose? Did Claude Code mention this?
    \item \textbf{Survey weights:} Did it use them unless you told it to?
    \item \textbf{Income = 0:} ``No income'' is coded as 0. How did Claude Code treat it?
    \item \textbf{Tests:} Did it copy the pipeline's own output in as an expected value, or edit one to get a pass?
  \end{itemize}

  \vspace{0.8em}
  \alert{These are the moments where your judgment matters.}\\
  The workflow is: notice the decision, direct Claude Code explicitly, then verify the output.
\end{frame}

\begin{frame}{Debrief}
  \textbf{Discussion questions:}

  \vspace{0.5em}
  \begin{enumerate}
    \small
    \item Did your merge lose any observations? How did you handle weekends?
    \item How many observations did you lose when you added controls? Did Claude Code tell you?
    \item How many of your robustness specifications were correct on the first try?
    \item What did you catch in the plan that you would have missed in the output?
    \item Did any test fail? Did Claude Code fix the code, or change the test?
  \end{enumerate}

  \vspace{0.5em}
  \begin{block}{Key takeaway}
    The tests, the assertions, and the plan review are not extra work --- they are the workflow. The 30 seconds it takes to read a plan or check a sample size can save hours of debugging downstream.
  \end{block}
\end{frame}
